\clearpage\section{Kanoniske transformationer og Poissonparenteser}
\subsection{Definition og Kroneckers delta}\label{sec:canon}
En lokalt invertibel transformation $(q,p)\mapsto(Q,P)$ er kanonisk, hvis de fundamentale Poissonparenteser bevares:
\key{\PB{Q_i}{Q_j}=0,\qquad\PB{P_i}{P_j}=0,\qquad\PB{Q_i}{P_j}=\delta_{ij}.}
Kroneckers delta er $\delta_{ij}=1$ for $i=j$ og $0$ for $i\ne j$. Ved den direkte test af $(Q,P)$ beregnes parenteserne i de \emph{gamle} variable $(q,p)$. For den inverse transformation kan man i stedet teste de gamle variable i de nye kanoniske variable.
\begin{note}\textbf{Ét koordinatpar:} $\{Q,Q\}=\{P,P\}=0$ automatisk; kontrollér kun $\{Q,P\}=Q_qP_p-Q_pP_q=1$.\end{note}
\begin{method}\textbf{Vælg det sæt, funktionerne er udtrykt ved.} Hvis $Q=Q(q,p)$ og $P=P(q,p)$ er opgivet, beregn
\[\PB QP_{q,p}=\pd Qq\pd Pp-\pd Qp\pd Pq=1.\]
Hvis der i stedet er opgivet $q=q(Q,P)$ og $p=p(Q,P)$, beregn
\[\PB qp_{Q,P}=\pd qQ\pd pP-\pd qP\pd pQ=1.\]
For \emph{én frihedsgrad} og en lokalt invertibel transformation er testene ækvivalente. For flere frihedsgrader skal alle de fundamentale parenteser testes, enten for den direkte eller for den inverse transformation. Bland ikke afledte med hensyn til forskellige koordinatsæt.\end{method}
\textbf{To koordinatpar:} Kontrollér de seks uafhængige relationer
\[\PB{Q_1}{P_1}=\PB{Q_2}{P_2}=1,\qquad\PB{Q_1}{P_2}=\PB{Q_2}{P_1}=0,\]
\[\PB{Q_1}{Q_2}=\PB{P_1}{P_2}=0.\]
Hver parentes rummer summen over $k=1,2$ af afledte i $(q_k,p_k)$.
\subsection{Poissonparenteser og regneregler}\label{sec:pb}
\key{\PB FG=\sum_{k=1}^n\left(\pd F{q_k}\pd G{p_k}-\pd F{p_k}\pd G{q_k}\right).}
\[\PB F{q_i}=-\pd F{p_i},\qquad\PB F{p_i}=\pd F{q_i},\qquad\PB FG=-\PB GF.\]
\begin{note}Eksempel fra rotationsgeneratoren, jf. afsnit~\ref{sec:angular}: $L_x=yp_z-zp_y$, derfor $\{L_x,y\}=-\partial L_x/\partial p_y=z$ og $\{L_x,p_y\}=\partial L_x/\partial y=p_z$.\end{note}
\subsection{Genererende funktioner: fire typer}
Alle fire typer antages at definere en lokalt invertibel kanonisk transformation. I en standardkonvention:
\begin{center}\begin{tabular}{llll}\toprule Type & Variable & Gammel impuls/koordinat & Ny impuls/koordinat\\\midrule
$F_1$&$(q,Q,t)$&$p=\partial_qF_1$&$P=-\partial_QF_1$\\
$F_2$&$(q,P,t)$&$p=\partial_qF_2$&$Q=\partial_PF_2$\\
$F_3$&$(p,Q,t)$&$q=-\partial_pF_3$&$P=-\partial_QF_3$\\
$F_4$&$(p,P,t)$&$q=-\partial_pF_4$&$Q=\partial_PF_4$\\\bottomrule\end{tabular}\end{center}
For alle typer er $K=H+\partial F_a/\partial t$ med den afledte taget ved fastholdte \emph{uafhængige argumenter til den valgte type}.
\subsubsection*{Genererende funktioner skrevet med Leibniz-notation}
\[\begin{aligned}
F_1(q,Q,t):\quad &p_i=\pd{F_1}{q_i}, &P_i&=-\pd{F_1}{Q_i},\\
F_2(q,P,t):\quad &p_i=\pd{F_2}{q_i}, &Q_i&=\pd{F_2}{P_i},\\
F_3(p,Q,t):\quad &q_i=-\pd{F_3}{p_i}, &P_i&=-\pd{F_3}{Q_i},\\
F_4(p,P,t):\quad &q_i=-\pd{F_4}{p_i}, &Q_i&=\pd{F_4}{P_i}.
\end{aligned}\]
\begin{method}\textbf{Kontrollér en given genererende funktion.} Find typen ud fra dens uafhængige argumenter; differentier med hensyn til begge argumenter; sammenlign de resulterende udtryk med den ønskede transformation. For eksempel kræver $S(q,P,t)$, at $p=\partial S/\partial q$ og $Q=\partial S/\partial P$. Tidsleddet $\partial S/\partial t$ bestemmer derimod ændringen af Hamiltonfunktionen og tester ikke alene transformationen.\end{method}
\begin{method}\textbf{Løs bevægelsen via en kanonisk transformation.} Omskriv $H(q,p,t)$ til $\widetilde H(Q,P,t)$, løs $\dot Q=\partial\widetilde H/\partial P$ og $\dot P=-\partial\widetilde H/\partial Q$, og indsæt derefter $Q(t),P(t)$ i $q(Q,P,t),p(Q,P,t)$.\end{method}
\subsection{Den transformerede Hamiltonfunktion}\label{sec:newH}
\begin{method}\textbf{1.} Identificér transformationen og den korrekte genererende funktion $F_a$. \textbf{2.} Beregn $\partial F_a/\partial t$ ved fastholdte argumenter. \textbf{3.} Invertér for $q=q(Q,P,t)$ og $p=p(Q,P,t)$. \textbf{4.} Indsæt begge i den oprindelige $H$ og tilføj tidsleddet.\end{method}
\key{\widetilde H(Q,P,t)=H(q,p,t)+\left.\pd{F_a}{t}\right|_{\text{egne argumenter}}.}
Her skal højresiden til sidst omskrives til $(Q,P,t)$. Hvis $F_a$ er \emph{eksplicit} tidsuafhængig, er tidsleddet nul, men $H$ skal stadig omskrives til $(Q,P)$.
\textbf{Mini-eksempel.} For $Q=aq$, $P=p/a$ med konstant $a\ne0$ og $H=p^2/(2m)+kq^2/2$ fås $q=Q/a$, $p=aP$, og
\[K(Q,P)=\frac{a^2P^2}{2m}+\frac{kQ^2}{2a^2}.\]
\begin{note}Tidsafhængige kanoniske transformationer kræver omhu: $\partial F/\partial t$ er ikke den totale afledte $\dd F/\dd t$.\end{note}


\subsubsection*{Omvendt opgavetype: find $F$ fra den ønskede $\widetilde H$}
\begin{method}\textbf{1.} Undersøg, om man kan vælge $\partial F/\partial t=0$ og dermed $\widetilde H=H$. Det er et muligt valg, ikke en almen regel. \textbf{2.} Brug den ønskede $\widetilde H$ til at vælge nye variable, fx $P=H$ hvis $\widetilde H=P$. \textbf{3.} Identificér $F_1$--$F_4$ og isolér den gamle variabel, der passer til dens argumenter. \textbf{4.} Integrér den tilhørende transformationsrelation og husk en vilkårlig funktion af de andre argumenter. \textbf{5.} Differentier for at bestemme den anden nye variabel, og kontrollér begge relationer samt $\widetilde H=H+F_t$.\end{method}
Eksempel: $H=p^2/(2m)+mgq$, ønsket $\widetilde H=P$ og en type 4-funktion $U(p,P)$. Vælg $U_t=0$ og $P=H$, så
\[q=\frac{P}{mg}-\frac{p^2}{2m^2g},\quad q=-U_p\quad\Rightarrow\quad\boxed{U(p,P)=-\frac{pP}{mg}+\frac{p^3}{6m^2g}+f(P)},\]
\[Q=U_P=-\frac{p}{mg}+f'(P).\]
Vælg eksempelvis $f=0$; $\{Q,P\}_{q,p}=1$. En direkte integration af $F_t=\widetilde H-H$ er generelt utilstrækkelig, fordi den ukendte transformation også bestemmer $q,p$ på højresiden.

\subsection{Oprindelse af de genererende funktioner}
En kanonisk transformation bevarer Hamiltons variationsprincip, hvis fase-rummets én-former opfylder
\[\sum_i p_i\,\dd q_i-H\,\dd t=\sum_i P_i\,\dd Q_i-K\,\dd t+\dd F_1(q,Q,t).\]
Sammenligning af koefficienterne foran $\dd q_i$, $\dd Q_i$ og $\dd t$ giver $p_i=\partial F_1/\partial q_i$, $P_i=-\partial F_1/\partial Q_i$ og $K=H+\partial F_1/\partial t$.
Ved Legendreomformninger fås $F_2=F_1+\sum_iP_iQ_i$, $F_3=F_1-\sum_i p_iq_i$ og $F_4=F_1+\sum_iP_iQ_i-\sum_i p_iq_i$; dette forklarer fortegnstabellen ovenfor.
\subsection{Symplektisk Jacobimatrix}
Sæt $z=(q,p)^T$, $Z=(Q,P)^T$, $M=\partial Z/\partial z$, $J=\left(\begin{smallmatrix}0&I\\-I&0\end{smallmatrix}\right)$. Den kompakte kanonicitetsbetingelse er
\[\boxed{MJM^T=J.}\]
Dette er ækvivalent med samtlige fundamentale Poissonparentesbetingelser. Under standardantagelser følger også $M^TJM=J$.

\subsection{Rotationsgeneratorer og bevarelse}
Den mekaniske rotationsgenerator i kanoniske kartesiske variable har $\vv L=\vv r\times\vv p$ for sædvanlige systemer; så $\PB{L_x}{y}=z$, $\PB{L_x}{z}=-y$, $\PB{L_x}{p_y}=p_z$. \emph{Generatorvirkningen} skal ikke forveksles med \emph{bevarelse}: for en tidsuafhængig rotationsinvariant $H$ er $\PB{L_x}{H}=0$, og dermed er $L_x$ bevaret.

